% !TeX program = xelatex
% 编译：TEXINPUTS='.:../../../.agents/skills/book-translation/templates/latex//:' xelatex ./2026-10-07-MirrorOPD-从PMD与BPO到访问测度.tex
\documentclass[lang=cn,scheme=chinese,thmcnt=chapter,device=normal]{elegantbook}

\usepackage{mathtools}
\usepackage{amssymb}
\usepackage{bm}
\usepackage{nicefrac}
\usepackage{booktabs}
\usepackage{enumitem}
\usepackage{float}

\newcommand{\E}{\mathbb{E}}
\newcommand{\Var}{\operatorname{Var}}
\newcommand{\Cov}{\operatorname{Cov}}
\newcommand{\KL}{\operatorname{KL}}
\newcommand{\Fisher}{\mathcal{F}}
\newcommand{\beh}{\mu}
\newcommand{\cand}{\pi}
\newcommand{\ratio}{\rho}
\newcommand{\stop}{\operatorname{stopgrad}}
\DeclareMathOperator*{\argmax}{arg\,max}
\DeclareMathOperator*{\argmin}{arg\,min}
\DeclareMathOperator{\supp}{supp}

\hypersetup{pdftitle={从 token-level 最优解统一推导 PMD、BPO 与 MDPO}}
\setcounter{tocdepth}{2}

\begin{document}
\frontmatter
\begin{center}
{\bfseries\LARGE 从 token-level 最优解统一推导 PMD、BPO 与 MDPO}\\[1.5ex]
{\large 梯度、重要性采样、偏差方差与实现比较}
\end{center}
\vspace{2ex}

\chapter*{摘要}
\addcontentsline{toc}{chapter}{摘要}
本文从固定 prefix 上的 KL 镜像子问题出发，统一写出 PMD-Part、PMD-Mean、BPO 与 MDPO 的优化目标、精确最优解和参数梯度。主线结论是：PMD-Part 的闭式 target 是 Gibbs 分布；PMD-Mean 是带单纯形约束的 Lambert-$W$ 投影；MDPO 直接优化 KL 正则化目标，梯度中保留 action importance ratio；BPO 从平方残差回归出发，在当前策略点附近线性化，再把梯度实现为带 mismatch weight 的 per-token policy-gradient surrogate。四者可以共享同一个 token-level 最优策略，但离开最优点后的梯度并不相同。重要性采样降低测度偏差却引入 ratio 方差，回归和 BPO 的截断权重降低方差却产生投影或截断偏差。最后核对 OpenKimi、经典 MDPO 和 BPO 论文源码，区分固定 inner budget 与真正的子问题收敛。

\tableofcontents
\mainmatter

\chapter{统一的 token-level 镜像子问题}
\label{ch:unified-problem}

\section{Prefix、行为策略与局部 reward}
\label{sec:notation}

固定一个 prefix $s=(x,a_{<t})$，词表动作为 $a$，行为策略为 $\beh(a\mid s)$，待更新学生策略为 $\cand_\theta(a\mid s)$。以下所有结论都在
\[
\cand_\theta(\cdot\mid s)\ll\beh(\cdot\mid s)
\]
的 support 条件下成立。记局部监督信号为 $f(s,a)$；它可以是即时 per-token reward、优势 $A^\beh(s,a)$，也可以是包含未来 reward 的 soft $Q$。

设 $d_\beh(s)$ 是行为轨迹访问 prefix 的测度。若一个条件 target 在每个 prefix 上可以独立实现，则 $d_\beh(s)>0$ 只决定该 prefix 在总体训练损失中的权重；如果模型共享参数、数据有限或存在未覆盖 prefix，它会改变投影点和估计方差。

\section{KL 镜像最优解}
\label{sec:token-optimum}

给定温度 $\tau>0$，token-level PMD 子问题为
\begin{equation}
\label{eq:token-subproblem}
 \cand^+(\cdot\mid s)=\argmax_{p\ll\beh(\cdot\mid s)}
 \left\{\E_{a\sim p}[f(s,a)]-\tau\KL(p\Vert\beh(\cdot\mid s))\right\}.
\end{equation}
令
\[
 Z_f(s)=\E_{a\sim\beh(\cdot\mid s)}\exp\left(\frac{f(s,a)}{\tau}\right).
\]

\begin{theorem}[Token-level Gibbs 解]
\label{thm:gibbs}
若 $0<Z_f(s)<\infty$，式 \eqref{eq:token-subproblem} 的唯一解为
\begin{equation}
\label{eq:gibbs}
 \cand^+(a\mid s)=\frac{\beh(a\mid s)\exp(f(s,a)/\tau)}{Z_f(s)}.
\end{equation}
等价地，$u^+(s,a)=\log(\cand^+/\beh)$ 满足
\begin{equation}
\label{eq:optimal-logratio}
 u^+(s,a)=\frac{f(s,a)}{\tau}-\log Z_f(s).
\end{equation}
\end{theorem}

\begin{proof}
由式 \eqref{eq:gibbs}，有
\[
\log\cand^+(a\mid s)=\log\beh(a\mid s)+f(s,a)/\tau-\log Z_f(s).
\]
代入 KL 得
\begin{align*}
\KL(p\Vert\cand^+)
&=\KL(p\Vert\beh)-\frac1\tau\E_p[f]+\log Z_f(s).
\end{align*}
整理得到
\[
\E_p[f]-\tau\KL(p\Vert\beh)
=\tau\log Z_f(s)-\tau\KL(p\Vert\cand^+).
\]
KL 非负且仅在 $p=\cand^+$ 时为零，故解唯一。
\end{proof}

\begin{corollary}[优势形式]
若 $f=A^\beh$ 且 $\E_{\beh}[A^\beh(s,\cdot)\mid s]=0$，则
\[
\cand^+(a\mid s)=\beh(a\mid s)\exp(A^\beh(s,a)/\tau) / Z_A(s).
\]
若 $A^\beh$ 是与序列 PMD 一致的 soft advantage，则 $Z_A(s)=1$；若它是硬 reward advantage，$Z_A(s)$ 一般不等于 $1$。
\end{corollary}

\section{即时 reward 与完整序列 reward}
\label{sec:immediate-vs-return}

若 $f(s,a)=r_t(s,a)$，式 \eqref{eq:token-subproblem} 是局部 token reward 的镜像更新。若完整目标是
\[
R(y)=\sum_{t=1}^{T}r_t(s_t,a_t),
\]
则与 sequence-level PMD 一致的 token score 必须是
\[
Q^\beh(s_t,a_t)=r_t(s_t,a_t)+\E_\beh\left[\sum_{k>t}r_k(s_k,a_k)\mid s_t,a_t\right],
\]
再取相应 soft advantage。直接把 $r_t$ 放进式 \eqref{eq:gibbs} 会得到局部目标，而不是完整序列目标。per-token reward 消除了终端奖励的稀疏性，却没有消除 future reward 对全局最优条件分布的作用。

\chapter{四种方法从最优解出发的推导}
\label{ch:four-methods}

\section{PMD-Part：回归最优 log-ratio}
\label{sec:pmd-part}

定义 $u_\theta(s,a)=\log\cand_\theta(a\mid s)/\beh(a\mid s)$，并令
\[
h(s,a)=f(s,a)/\tau-\log Z_f(s).
\]
PMD-Part 在行为测度下回归精确最优性方程：
\begin{equation}
\label{eq:part-loss}
 \mathcal L_{\mathrm{Part}}(\theta)
 =\frac12\E_{s\sim d_\beh,\,a\sim\beh}
 \left[u_\theta(s,a)-h(s,a)\right]^2.
\end{equation}
把 target $h$ 和行为策略视为固定，记 $g_\theta(s,a)=\nabla_\theta\log\cand_\theta(a\mid s)$，则
\begin{equation}
\label{eq:part-gradient}
 -\nabla_\theta\mathcal L_{\mathrm{Part}}
 =\E_{d_\beh\beh}\left[(h-u_\theta)g_\theta\right].
\end{equation}

\begin{proposition}[PMD-Part 的零点]
\label{prop:part-zero}
在无限容量、正访问权重和满 action support 下，$\mathcal L_{\mathrm{Part}}=0$ 当且仅当 $\cand_\theta(\cdot\mid s)=\cand^+(\cdot\mid s)$ 对 $d_\beh$-几乎处处的 prefix 成立。
\end{proposition}

\begin{proof}
平方损失为非负数。它为零当且仅当 $u_\theta=h$ 在行为测度下几乎处处成立。指数化并使用式 \eqref{eq:optimal-logratio} 即得 $\cand_\theta=\cand^+$。反向蕴含显然成立。
\end{proof}

PMD-Part 的重要性质是目标中没有 ratio 权重；ratio 只出现在被拟合的 log-ratio 中。因此它避免了直接 policy gradient 的乘法重尾，但有限模型下它优化的是 $d_\beh$ 加权的投影。

\section{PMD-Mean：中心化回归与 Lambert-$W$}
\label{sec:pmd-mean}

令
\[
\bar f(s)=\E_{a\sim\beh}[f(s,a)\mid s],
\qquad g(s,a)=\frac{f(s,a)-\bar f(s)}{\tau}.
\]
PMD-Mean 不显式估计 $\log Z_f(s)$，而求解
\begin{equation}
\label{eq:mean-loss}
 \min_u\;\frac12\E_{s\sim d_\beh,a\sim\beh}[(u-g)^2]
 \quad\text{s.t.}\quad
 \E_{a\sim\beh(\cdot\mid s)}[e^{u(s,a)}]=1\quad\forall s.
\end{equation}
约束是策略归一化，而不是可省略的附加条件。

\begin{theorem}[PMD-Mean 的逐 prefix 解]
\label{thm:mean-lambert}
在有限 action space、$\beh(a\mid s)>0$ 和可行条件下，式 \eqref{eq:mean-loss} 的 KKT 解满足
\begin{equation}
\label{eq:mean-lambert}
 u^+(s,a)=g(s,a)-W\!\left(\lambda(s)e^{g(s,a)}\right),
\end{equation}
其中 $\lambda(s)\geq0$ 由 $\E_\beh[e^{u^+(s,\cdot)}\mid s]=1$ 唯一确定。
\end{theorem}

\begin{proof}
对固定 $s$ 加入约束乘子，拉格朗日函数为
\[
\frac12\E_\beh[(u-g)^2]+\lambda(s)(\E_\beh[e^u]-1).
\]
逐 action 的一阶条件是 $u-g+\lambda e^u=0$。令 $v=g-u$，则 $ve^v=\lambda e^g$，所以 $v=W(\lambda e^g)$，得到式 \eqref{eq:mean-lambert}。归一化条件确定 $\lambda$；当 $\lambda\geq0$ 时左侧随 $\lambda$ 单调下降，故解唯一。
\end{proof}

PMD-Part 的修正项 $\log Z_f(s)$ 对 action 是常数，PMD-Mean 的 Lambert-$W$ 项依赖 action。两者只有在退化或额外近似下才相同；它们可以拥有相同的排序和单调改进性质，但不是同一个 mirror subproblem 的两个无偏实现。

\section{MDPO：直接优化 KL 正则化目标}
\label{sec:mdpo}

MDPO 直接对式 \eqref{eq:token-subproblem} 的目标做参数优化。令
\[
F_s(\theta)=\E_{a\sim\cand_\theta}[f(s,a)]-\tau\KL(\cand_\theta\Vert\beh).
\]
对 $\theta$ 求导，利用 $\E_{\cand}[g_\theta]=0$，得到
\begin{equation}
\label{eq:mdpo-onpolicy-gradient}
\nabla F_s
=\E_{a\sim\cand}\left[
\left(f(s,a)-\tau\log\ratio_\theta(s,a)\right)g_\theta(s,a)
\right],
\end{equation}
其中 $\ratio_\theta=\cand_\theta/\beh$。若样本来自行为策略，则
\begin{equation}
\label{eq:mdpo-is-gradient}
\nabla F_s
=\E_{a\sim\beh}\left[
\ratio_\theta(s,a)\left(f(s,a)-\tau\log\ratio_\theta(s,a)\right)g_\theta(s,a)
\right].
\end{equation}
这就是 MDPO 中的 action-level IS。若状态也来自 replay 的 $d_\beh$，而目标是 $d_{\cand}$ 下的 occupancy 目标，还需要 prefix ratio；经典实现通常把状态测度固定为 replay/rollout 测度，因而没有完整解决 state-distribution mismatch。

\begin{proposition}[MDPO 梯度与 PMD-Part 回归梯度的零点一致、方向不一致]
\label{prop:mdpo-part}
MDPO 的梯度式 \eqref{eq:mdpo-is-gradient} 与 PMD-Part 梯度式 \eqref{eq:part-gradient} 在精确镜像解处都为零，但一般不相等。
\end{proposition}

\begin{proof}
在镜像解处，$\log\ratio=f/\tau-\log Z_f$，故式 \eqref{eq:mdpo-onpolicy-gradient} 为最优目标的驻点；PMD-Part 的 residual $h-u$ 逐点为零。离开镜像解时，MDPO 被积函数包含 $\ratio_\theta$ 和 $f-\tau\log\ratio_\theta$，PMD-Part 被积函数只有 $h-u_\theta$，两者测度和残差均不同。
\end{proof}

\section{BPO：从 token 最优性到实际 PG surrogate}
\label{sec:bpo}

BPO 的 token-level PMD 起点与式 \eqref{eq:token-subproblem} 相同，使用 $f=A^\beh$。在终端奖励情形，论文先定义 token residual
\begin{equation}
\label{eq:bpo-token-residual}
 d_\pi(s_t,a_t)=\eta A^\beh(s_t,a_t)-\left[\log\frac{\cand(a_t\mid s_t)}{\beh(a_t\mid s_t)}+\KL(\beh(\cdot\mid s_t)\Vert\cand(\cdot\mid s_t))\right].
\end{equation}
该 residual 在 $a_t\sim\beh(\cdot\mid s_t)$ 下条件均值为零。BPO 将沿轨迹的 residual 累加成平方回归目标，并证明其零点与 PMD 解一致。

实际算法没有把式 \eqref{eq:bpo-token-residual} 的平方损失做 inner convergence。论文在 $\cand=\beh$ 处线性化，使用组内 reward mean/std 估计初始 value 和缩放因子，再以二元 KL 近似逆向 KL，得到
\begin{equation}
\label{eq:bpo-gradient}
 \nabla\mathcal L_{i,t}^{\mathrm{BPO}}
 \approx -A_i M_{i,t}\min(\widetilde w_{i,t},C)
 \nabla\log\cand(a_{i,t}\mid s_{i,t}),
\end{equation}
其中
\begin{equation}
\label{eq:bpo-weight}
 \widetilde w_{i,t}=rac{1+\epsilon-\beh(a_{i,t}\mid s_{i,t})}{1+\epsilon-\cand(a_{i,t}\mid s_{i,t})}.
\end{equation}
因此 BPO 最终训练的是 PG 形式的 surrogate；回归是理论构造和梯度来源，而不是独立的第二阶段优化器。

\chapter{梯度的偏差与方差}
\label{ch:bias-variance}

\section{三个梯度估计器}
\label{sec:estimators}

固定 prefix，设 $g=\nabla\log\cand_\theta(a\mid s)$，$\rho=\cand_\theta/\beh$，并定义 $c=f-\tau\log\rho$。MDPO 的无偏行为测度估计器为
\[
\widehat G_{\mathrm{MDPO}}=\rho c g,
\qquad
\E_\beh[\widehat G_{\mathrm{MDPO}}]=\nabla F_s.
\]
PMD-Part 的回归估计器为
\[
\widehat G_{\mathrm{Part}}=(h-u_\theta)g,
\]
它在 target 固定且回归残差正确时收敛到镜像解，但一般不是 $\nabla F_s$ 的无偏估计。BPO 的实际估计器为
\[
\widehat G_{\mathrm{BPO}}=-A M\min(\widetilde w,C)g.
\]
由于 $M$ 和截断依赖当前 ratio，它也一般是 surrogate 的有偏估计器。

\begin{theorem}[IS 梯度的二阶矩上界]
\label{thm:is-variance}
若 $\|g\|\leq G$ 且 $|c|\leq C_f$，则
\begin{equation}
\label{eq:is-second-moment}
 \E_\beh[\|\widehat G_{\mathrm{MDPO}}\|^2]
 \leq G^2C_f^2\E_\beh[\rho^2].
\end{equation}
因此其方差满足 $\Var(\widehat G_{\mathrm{MDPO}})\leq G^2C_f^2\E_\beh[\rho^2]$；当 ratio 重尾或 support 接近失败时，该上界迅速变大。
\end{theorem}

\begin{proof}
逐样本有 $\|\rho c g\|^2\leq \rho^2C_f^2G^2$。取行为期望并使用 $\Var(X)\leq\E[\|X\|^2]$ 即得。
\end{proof}

\begin{proposition}[有界 mismatch weight 的方差与截断偏差]
\label{prop:bpo-tradeoff}
若 $0\leq M\leq1$ 且 $0\leq\min(\widetilde w,C)\leq C$，$|A|\leq C_A$，$\|g\|\leq G$，则
\[
\E[\|\widehat G_{\mathrm{BPO}}\|^2]
\leq C_A^2C^2G^2.
\]
相对于未截断的 surrogate 梯度 $G_0=-AM\widetilde w g$，截断引入的 bias 满足
\[
\left\|\E[\widehat G_{\mathrm{BPO}}-G_0]\right\|
\leq C_A G\E[(\widetilde w-C)_+].
\]
\end{proposition}

\begin{proof}
第一式由逐样本范数上界直接得到。第二式使用三角不等式以及 $|\min(\widetilde w,C)-\widetilde w|=(\widetilde w-C)_+$。
\end{proof}

\begin{remark}[偏差方差位置]
MDPO 的 action IS 主要承担测度校正，理想情况下无偏但方差可能很大。PMD-Part 去掉了梯度外部 ratio，代价是离开镜像根时不再等于原始 KL 目标梯度，并且受到 $d_\beh$ 加权投影影响。PMD-Mean 额外用均值替代 log-partition，产生几何投影偏差。BPO 的 mismatch weight 比标准 ratio 更温和，并通过平滑、mask 和截断控制方差，但因此不再是原始 PMD 梯度的无偏估计。
\end{remark}

\section{prefix IS 的位置}
\label{sec:prefix-is}

若数据来自 $d_\beh(s)\beh(a\mid s)$，而要无偏估计目标 occupancy 下的量，则完整权重为
\[
\frac{d_\cand(s)}{d_\beh(s)}\frac{\cand(a\mid s)}{\beh(a\mid s)}=W_{<t}\rho_t.
\]
PMD-Part、PMD-Mean 和 BPO 的条件回归形式通常只使用 $d_\beh$ 访问的 prefix，因此没有显式 $W_{<t}$；这不是一般 off-policy 无偏性，而是把目标定义为行为测度下的 pointwise regression。若在有限容量下用 prefix IS 改变投影权重，必须同时承担 $W_{<t}$ 的方差和 support 风险。

\chapter{实现与参数比较}
\label{ch:implementation}

\section{理论对象与实际代码}
\label{sec:implementation-table}

\begin{table}[H]
\centering
\caption{PMD、MDPO 与 BPO 的目标和实现差异}
\label{tab:method-comparison}
\begin{tabular}{p{0.16\textwidth}p{0.23\textwidth}p{0.23\textwidth}p{0.25\textwidth}}
\toprule
方法 & 理论更新 & 梯度/权重 & 主要偏差方差来源 \\
\midrule
PMD-Part & 回归 $u^+=f/\tau-\log Z_f$ & 行为测度下 residual score，无外部 IS & 有限模型投影偏差；通常低于 IS 的 ratio 方差 \\
PMD-Mean & 中心化平方回归加 $\E_\beh[e^u]=1$ & Lambert-$W$ 投影梯度 & 均值和投影近似偏差；无乘积 IS \\
MDPO & 直接最大化 $\E_\pi[f]-\tau\KL(\pi\Vert\mu)$ & $\rho(f-\tau\log\rho)\nabla\log\pi$ & action IS 无偏但二阶矩随 $\E_\mu[\rho^2]$ 增长 \\
BPO & residual regression 的局部线性化 & per-token $A\,M\min(\widetilde w,C)\nabla\log\pi$ & 线性化、二元 KL、截断产生 bias，权重受控 \\
\bottomrule
\end{tabular}
\end{table}

\section{BPO 论文的训练预算}
\label{sec:bpo-implementation}

BPO 论文源码的主实验写明：每个 training step 先生成 256 个 prompt、每个 prompt 16 个 response，共 4096 条 response；再划分为 8 个、每个 512 条 response 的 mini-batch，执行 8 次 optimizer updates。主实验运行 400 个 training steps，即 3200 次 optimizer updates；附录消融运行 1000 个 training steps。论文没有按 residual norm 或 mirror-subproblem gap 停止的 inner loop，因此这些是固定更新预算，不是子问题收敛步数。

\section{OpenKimi/PMD 的训练预算}
\label{sec:openkimi-implementation}

OpenKimi 的 PMD loss 默认先对完整 response 的 token log-ratio 求和，再把 sequence log-ratio broadcast 到 response 内 token。Qwen2.5-7B math recipe 设置 `PPO\_EPOCHS=1`、`train\_batch\_size=512`、`actor.ppo\_mini\_batch\_size=32`，所以静态 batch 下每个 outer step 约有 $512/32=16$ 个 mini-batch 更新。代码没有 inner residual stopping 或 PMD gap 检验；`total\_training\_steps=3000` 是 outer steps 上限。

\section{经典 MDPO 的训练预算与 IS}
\label{sec:mdpo-implementation}

经典 MDPO on-policy 代码默认 `sgd\_steps=5`，MuJoCo README 的推荐示例使用 10。一个 rollout batch 固定后，对同一批数据重复执行该数量的 policy optimizer updates，再刷新 old/new anchor；没有根据 mirror residual 自适应停止。

off-policy 代码的 README 暴露 `sgd\_steps=1000`，但当前训练循环实际写为 `for grad\_step in range(1)`，原本的 `int(self.gradient\_steps)` 被注释。因此每个 train event 实际只执行一次 actor/critic update；`gradient\_steps` 主要影响 target/anchor policy 同步周期。实现使用 replay 状态测度和 action-level policy gradient，不能解释为完整的 prefix-IS 求解。

\section{比较时应报告的指标}
\label{sec:reporting}

为了比较四种方法，实验至少应固定同一 rollout batch，并记录：每个 inner update 的 mirror loss、log-ratio residual norm、$\KL(\cand_\theta\Vert\beh)$、梯度范数、ratio 的均值与二阶矩、BPO mismatch weight 的截断比例，以及 prefix coverage。这样才能区分“做了更多 optimizer steps”和“真正接近 token-level mirror optimum”。

\chapter{结论与 MirrorOPD 的边界}
\label{ch:conclusion}

从同一个 token-level Gibbs 最优解出发，PMD-Part 是精确 log-ratio 回归，PMD-Mean 是中心化回归后的 Lambert-$W$ 投影，MDPO 是带 action IS 的直接 KL 正则化 policy gradient，BPO 是平方 residual 回归经过局部线性化和二元 KL 化简后的 per-token PG surrogate。它们在理想条件下可以共享相同的零点，但训练过程中的梯度、访问测度、偏差和方差不同。

本文件暂不推导 OPD 的完整 off-policy correction 或 KL 系数调度。后续 MirrorOPD 需要先选择局部即时 reward 目标还是完整加性回报目标，再决定是否引入 prefix IS、是否采用 Part/Mean 投影，以及如何用 Natural Gradient 控制学生策略与教师策略的距离。

\backmatter
\begin{thebibliography}{99}
\bibitem{tomar2020mdpo}Tomar et al. \emph{Mirror Descent Policy Optimization}. arXiv:2005.09814, 2020.
\bibitem{xu2026pmd}Xu et al. \emph{Approximation of Log-Partition Function in Policy Mirror Descent Induces Implicit Regularization for LLM Post-Training}. arXiv:2602.05933, 2026.
\bibitem{song2026bpo}Song et al. \emph{Bellman Policy Optimization}. arXiv:2609.15987, 2026.
\bibitem{openkimi}OpenKimi contributors. \emph{OpenKimi: PMD and RL post-training implementation}. GitHub repository, 2026. \url{https://github.com/horizon-llm/OpenKimi}.
\bibitem{mdpo-repo}Tomar et al. \emph{Mirror Descent Policy Optimization}: official implementation. GitHub repository, 2020. \url{https://github.com/manantomar/Mirror-Descent-Policy-Optimization}.
\end{thebibliography}
\end{document}
